\chapter{What an Energy Balance Unit is}\label{ch:action}

We can now define the account precisely. EBU measures the finite change of a declared physical potential. The definition is simple enough to calculate for one transfer, and strong enough to support the rest of the book: path accounting, complete transaction histories, joint action, feedback and equilibrium probability all connect back to this same finite comparison.


\section{A physical consequence needs an account}
The previous chapters have brought us to a practical question. The clinic cannot operate without payments, but a successful payment is not evidence that the water system remains capable of serving it. Conversely, water can be physically available while the people who need it cannot gain access. We do not have to abolish one description to improve the other. We need a second account that makes physical consequences visible at the point of action.

An Energy Balance Unit, or EBU, is the signed unit of that account. Given a declared physical potential $V$, an accepted finite action has the decrease of $V$ between its complete before-state and after-state. If the represented state moves toward the declared reference, the EBU can be positive; if it moves away, it can be negative. Zero means that this particular potential did not change. It does not mean that no water moved, no energy was used, or no human work occurred.

This definition is deliberately connected to a physical event. A promise to deliver water is not the delivery. A payment for it is not the delivery. The account must identify the state before the action, the change that actually occurred, and the physical boundary within which that change is assessed. It is equally important to name what the account is measuring: a field built from specified carriers and functions, not a moral score for the person who acted.

\section{The finite definition comes first}
Let $x^{\mathrm{b}}$ be the complete represented state before an action and let $x^{\mathrm{b}}+\delta$ be its complete successor under one fixed potential. Define
\begin{equation}\label{eq:intro-ebu-endpoint}
 E=V(x^{\mathrm{b}})-V(x^{\mathrm{b}}+\delta).
\end{equation}
The minus sign makes a decrease in represented burden positive. No gradient approximation, motion law or actor payment is needed to define this finite comparison. The field, physical boundary and complete endpoints must be declared before the number has meaning.

\section{Following the change along a path}
Recall the bent ray from Chapter~\ref{ch:research}. Fermat's optical travel-time integral can depend on the spatial route. A fixed-potential EBU integral has a different structure: it integrates the differential of a supplied state function. Its role is to explain how changing local marginals add to the finite endpoint value already defined.

Choose an admissible piecewise $C^1$ path $\gamma$ from $x^{\mathrm{b}}$ to $x^{\mathrm{b}}+\delta$, and suppose $V$ is $C^1$ on a neighbourhood of that path. The chain rule and fundamental theorem of calculus give
\begin{equation}\label{eq:intro-ebu-path}
 E=-\int_0^1\nabla V(\gamma(s))^{\mathsf T}\gamma'(s)\,ds.
\end{equation}
Chapter~\ref{ch:path} proves this identity carefully. A straight evaluation path need not be an observed physical trajectory; that interpretation requires a physical response model. The total is path independent under these assumptions because it is the difference of one single-valued potential.

\begin{intuitionbox}{One definition, a conditional path representation}
Finite EBU is the before-minus-after potential difference. Under the regularity assumptions, a changing-field integral gives the same number. Neither expression grants physical permission, chooses a motion law or assigns ownership.
\end{intuitionbox}

\section{Where the rest of the physical process belongs}
A water transfer may require pumping, filtration or another physical transformation. Those consequences do not become invisible merely because our first example follows water alone. If the pump consumes electricity, a fuller physical model records the energy carrier and its transformation. If water is lost, the lost amount enters a declared sink coordinate or crosses an explicit boundary. A potential defined on that fuller state can then evaluate all the consequences it is meant to value.

This is different from subtracting an unexplained ``process cost'' from the EBU equation. The general finite definition in this book remains
\begin{equation}\label{eq:intro-ebu-net}
 E=V(x^{\mathrm{b}})-V(x^{\mathrm{b}}+\delta).
\end{equation}
An engineering or institutional account may keep additional, separately named quantities. Such a quantity is not physical EBU by a change of notation, and the same effect must not be counted twice. Our first water transfer is an intentionally lossless teaching world. Its simplicity lets us learn the equation; it does not assert that operating a real pump is free.

\section{A four-unit transfer}
A holds $18\X$, B holds $2\X$, and the declared reference is $(10,10)$. Here $\X$ denotes one unit of the represented physical stock. Each location has scale $2\X$. The potential we will construct in Chapter~\ref{ch:potential} is
\begin{equation}
 V(x)=\frac12\left(\frac{x_1-10}{2}\right)^2+
      \frac12\left(\frac{x_2-10}{2}\right)^2 .
\end{equation}
Its value at $(18,2)$ is $16$. If four units are transferred from A to B, the state becomes $(14,6)$ and the potential becomes $4$. The exact EBU is therefore $12\E$. Four physical units arrived at B; twelve EBU express the decrease of this particular potential. They are different units describing different properties of the same event.

If eight units are transferred, the endpoint is the reference and the EBU is $16\E$. Transferring sixteen units instead reaches $(2,18)$, with potential $16$ once again, so its EBU is zero. The last movement was not empty: it simply overshot far enough to cancel the initial gain in this account. At seventeen units the final potential is $20.25$ and the EBU is $-4.25\E$. These calculations show why a fixed reward for each delivered stock unit cannot substitute for valuation along the full finite action.

\diagram{transfer}{A four-unit lossless transfer preserves twenty physical units while the declared potential falls from sixteen to four.}{fig:transfer}

\section{What would an actor receive EBU for?}
Imagine two feasible actions from the same measured state. One restores a depleted supply without pushing another represented location into trouble. The other earns the same sales revenue while worsening the declared balance. A conventional invoice can look alike in both cases. The EBU account distinguishes their represented physical consequences. That distinction is the proposed incentive: useful restoration becomes visible and potentially rewardable, while depletion cannot be made to look like restoration merely by earning more money.

There is an important institutional step between a positive EBU calculation and a spendable personal balance. An institution would have to specify who may claim the action, how multiple contributors are treated, when a measured outcome becomes a receipt, and which future actions such a receipt could enable. A capacity design might allow demonstrated restorative value to support later physically permissible activity. The finite definition supplies a verifiable total for such a design. Attribution and the rules for using a receipt supply the remaining institutional structure. Later chapters keep a group result, an individual attribution and an institutional rule distinct.

The incentive also cannot be reduced to making every EBU number as large as possible. A well-functioning system may need continuing service and repair even when the represented state is already near its reference. An action may be essential despite a negative result in one field, and access to a clinic must not be made contingent on a person's ability to manufacture credits. The physical account can show a cost or benefit honestly; institutions must decide how essential provision, rights and obligations are protected. The resulting design can recognise maintenance, protect essential access and account honestly for resource use at the same time. These functions are compatible because the physical valuation is not made to carry every institutional decision.

\section{Permission and value are different questions}
Before a transfer is accepted, the source must have enough usable stock, the route must exist and the declared physical constraints must hold. A positive hypothetical quote for an impossible transfer creates no water. A negative quote for an urgent, permitted service does not by itself forbid the service. Permission answers whether the physical change is allowed; EBU answers how the declared potential changes.

Likewise, neither the source's loss nor the destination's gain alone determines the group result when several actions occur together. The potential is evaluated for the actual combined state change. Its exact group total can be known without a justified division into personal claims. We will keep that distinction as the mathematics develops.

Several useful properties follow from this finite definition. Successive receipts under one field telescope to the net endpoint change. Joint actions can be compared against all their subgroups to identify exactly what coordination contributes. A response model can supply the endpoint at a stated horizon, including pending flows and feedback. These are complementary capabilities of one account: reliable history, visible joint dependence and a clear connection to physical dynamics. We have named the unit and seen its consequence in one small world. The next chapter draws the world as cells and routes, so that a sentence such as four units reached B becomes an action we can test for physical possibility before calculating its value.
